\section{Modular localization and averaging}\label{sec:modular}

The recovery argument requires two kinds of control: modular spectral
localization must imply a predual norm estimate, and modular averaging
must control the moments of an operator even when its norm is large.
We establish these facts here. The two norm limits in
\Cref{mod:means} will control different terms in the fourth-moment
expansion in \Cref{sec:recovery}.

Throughout this Section $M$ is a von Neumann algebra with a faithful
normal state $\phi$. Let $H=L^2(M)$ be its standard Hilbert space,
$\xi=\phi^{1/2}$ its cyclic separating unit vector, and
\[
  L(a)h=ah,\qquad R(a)h=ha,\qquad
  D=\log\Delta_\phi,\qquad U_t=e^{itD}.
\]
Write $E_D$ for the spectral resolution of $D$ and $J$ for the modular
conjugation. Inner products are linear in the first variable.
We use the standard modular theory of \cite{Takesaki2003}.
We write $\sigma_t=\sigma^\phi_t$.
Standard modular theory gives
\begin{equation}\label{eq:modular-identities}
  U_t(a\xi)=\sigma_t(a)\xi,\qquad
  e^{D/2}a\xi=\xi a,\qquad
  R(a)=JL(a)^*J.
\end{equation}
Indeed, $M\xi$ is contained in the domain of the closed Tomita operator
$S=J\Delta_\phi^{1/2}$, with $S(a\xi)=a^*\xi$.
Thus $\Delta_\phi^{1/2}a\xi=Ja^*\xi=\xi a$.
These standard-form conventions are also those of
\cite[Section~2.1]{AHHM2020}.

\begin{lemma}[From modular localization to predual norm]\label{mod:predual}
For $w\in M$ and $s\in\R$,
\begin{equation}\label{eq:predual-bound}
  \norm{\phi w-e^sw\phi}_{M_*}
  \le (1+e^{s/2})\norm{(e^{D/2}-e^{s/2})w\xi}.
\end{equation}
Consequently, if $(w_n)$ is uniformly bounded and the $D$-spectral
support of $w_n\xi$ is contained in
$[s-d_n,s+d_n]$ with $d_n\to0$, then the left side of
\eqref{eq:predual-bound} tends to zero.
\end{lemma}

\begin{proof}
Put $c=e^{s/2}$ and $\eta=\xi w-cw\xi$.
By \eqref{eq:modular-identities},
$\eta=(e^{D/2}-c)w\xi$, and
$J\eta=w^*\xi-c\xi w^*$.
For $\norm{y}\le1$, commutation of the left and right actions gives
\[
  \langle y\xi,\xi w^*\rangle
  =\langle y(\xi w),\xi\rangle.
\]
Starting from $\phi(wy)=\langle y\xi,w^*\xi\rangle$, replace
$w^*\xi$ by $c\xi w^*$, and then $\xi w$ by $cw\xi$.
The total error is at most $(1+c)\norm{\eta}$, while the resulting
expression is $c^2\phi(yw)$. Taking the supremum over $y$ proves the
estimate. The last assertion follows from
\[
  \norm{(e^{D/2}-e^{s/2})w_n\xi}
  \le \sup_{\abs{u-s}\le d_n}\abs{e^{u/2}-e^{s/2}}\,
       \norm{w_n\xi}.
  \qedhere
\]
\end{proof}

\begin{lemma}[Smooth spectral localization]\label{mod:localization}
Let $f\in C_c^\infty(\R)$, with Fourier convention
$f(u)=\int_{\R}g(t)e^{itu}\,dt$.
For $y\in M$, the weak integral
\[
  y_f=\int_{\R}g(t)\sigma_t(y)\,dt
\]
belongs to $M$ and satisfies
$\norm{y_f}\le\norm{g}_1\norm{y}$ and
$y_f\xi=f(D)y\xi$.
If $a\in M$ and $a\xi$ has compact $D$-spectral support, there is $a^\flat\in M$
such that
\begin{equation}\label{eq:right-multiplier}
  a\xi=\xi a^\flat.
\end{equation}
The same Fourier construction applies to any continuous
state-preserving action and its unitary implementation.
\end{lemma}

\begin{proof}
The inverse Fourier transform $g$ is integrable.
The bounded weak integral is defined by integration against normal
functionals, and its operator norm has the stated bound.
Applying it to $\xi$ and using the spectral theorem gives
$y_f\xi=f(D)y\xi$. This uses no operator norm continuity of the orbit.

If $K$ is a compact support for $a\xi$, choose $f\in C_c^\infty(\R)$
equal to $e^{-u/2}$ on a neighborhood of $K$, and put $a^\flat=a_f$.
Then \eqref{eq:modular-identities} and functional calculus give
\[
  \xi a^\flat=e^{D/2}a^\flat\xi
  =e^{D/2}f(D)a\xi=a\xi.
\]
All vectors in this computation belong to the domain of $e^{D/2}$;
in particular the filtered vector has compact spectral support.
The Fourier argument for another state-preserving action is identical.
\end{proof}

For the rest of this Section assume $M_\phi=\C1$.
Let $\mu_k$ be uniform probability measure on $[-k,k]$, for
$k\in\N$.

\begin{lemma}[Two norm ergodic limits]\label{mod:means}
For every $y\in M$ and $\psi\in M_*$,
\begin{align}
  \int\sigma_t(y)\xi\,d\mu_k(t)&\longrightarrow\phi(y)\xi
       &&\text{in }H,\label{eq:hilbert-mean}\\
  \int\psi\circ\sigma_{-t}\,d\mu_k(t)&\longrightarrow\psi(1)\phi
       &&\text{in }M_*.\label{eq:predual-mean}
\end{align}
\end{lemma}

\begin{proof}
The unitary mean ergodic theorem gives a norm limit for the vectors in
\eqref{eq:hilbert-mean}. The corresponding operator averages have norm
at most $\norm{y}$. For each fixed $u$, translating the interval changes
their averages in operator norm by at most
$\abs{u}\norm{y}/k$ for all sufficiently large $k$.
Every weak operator cluster point is consequently fixed by $\sigma$,
hence scalar. Evaluation under $\phi$ identifies it as $\phi(y)1$.
This identifies the Hilbert norm limit and proves
\eqref{eq:hilbert-mean}.

The action on $M_*$ is norm continuous on each functional.
For positive normal functionals this follows from their realization as
vector functionals in standard form and the strong continuity of $U$;
linear decomposition gives the general case.
The predual averages are therefore Bochner integrals and are
contractions. Let
\[
  \mathcal N=\overline{\linspan}^{\,\norm{\cdot}}
   \{\rho\circ\sigma_u-\rho:\rho\in M_*,\ u\in\R\}.
\]
Its annihilator in $(M_*)^*=M$ is $M_\phi=\C1$.
The Hahn--Banach theorem implies
$\mathcal N=\{\psi\in M_*:\psi(1)=0\}$.
The averages tend to zero in norm on each generator of $\mathcal N$,
by the interval boundary estimate, and hence on all of $\mathcal N$ by
contractivity. They fix $\phi$. Apply this to $\psi-\psi(1)\phi$ to
obtain \eqref{eq:predual-mean}.
\end{proof}

Fix a free ultrafilter $\omega$ on $\N$. For a bounded continuous scalar
function $f$ define
\[
  m_t f(t)=\lim_{k\to\omega}\int f(t)\,d\mu_k(t).
\]
This positive normalized mean is invariant under translations and under
$t\mapsto-t$. For a bounded continuous $H$-valued function $F$, put
\[
  \norm{F}_{\mathrm{av}}
  =\bigl(m_t\norm{F(t)}^2\bigr)^{1/2}.
\]
The triangle inequality follows from the corresponding inequalities for
the finite interval integrals. Thus this is a seminorm.

\Needspace{6\baselineskip}
\begin{lemma}[Operator averaging]\label{mod:average}
For $S\in\B(H)$, define $P(S)$ by
\[
  \langle P(S)h,k\rangle
  =m_t\langle U_t^*SU_th,k\rangle.
\]
Then $P$ is a positive unital contraction, $P(S)$ commutes with every
spectral projection of $D$, and
\begin{equation}\label{eq:scalar-average}
  P(L(a))=\phi(a)1\qquad(a\in M).
\end{equation}
In particular,
\begin{equation}\label{eq:quadratic-average}
  \norm{t\mapsto SU_th}_{\mathrm{av}}^2
  =\langle P(S^*S)h,h\rangle.
\end{equation}
\end{lemma}

\begin{proof}
The bounded sesquilinear form defines $P(S)$, with
$\norm{P(S)}\le\norm{S}$. Positivity and unitality follow directly from
the mean. Translation invariance gives $U_u^*P(S)U_u=P(S)$ for every
$u$, hence commutation with the spectral projections of $D$.
For $S=L(a)$ the averages and their weak limits belong to the weakly
closed algebra $L(M)$, since
$U_t^*L(a)U_t=L(\sigma_{-t}(a))$.
The fixed-point algebra is scalar, and evaluation at $\xi$ gives
\eqref{eq:scalar-average}. Equation \eqref{eq:quadratic-average} is the
definition applied to $S^*S$.
No normality of $P$ is asserted or needed.
\end{proof}
